\documentclass[11pt]{article}

\usepackage{amsmath,amssymb,amsthm}
\usepackage{graphicx}
\usepackage{hyperref}
\usepackage{booktabs}

\newtheorem{theorem}{Theorem}
\newtheorem{conjecture}{Conjecture}
\newtheorem{proposition}{Proposition}
\newtheorem{remark}{Remark}

\title{On the Variance of the Volume Deficit of Random Inscribed
Polytopes in~$S^2$}

\author{Igor Rivin}

\date{\today}

\begin{document}

\maketitle

\begin{abstract}
Let $P_N$ denote the convex hull of $N$ independent uniformly distributed
points on the unit sphere $S^2 \subset \mathbb{R}^3$.  It is known
(Buchta--M\"uller--Tichy, Akopyan--Edelsbrunner--Nikitenko) that the expected
volume satisfies the exact formula
$\mathbb{E}[\operatorname{vol}(P_N)] = \frac{4\pi}{3}\,
\frac{(N-1)(N-2)(N-3)}{(N+1)(N+2)(N+3)}$,
which gives $\mathbb{E}[\Delta_N] = 16\pi/N + O(1/N^2)$ for the volume
deficit $\Delta_N = \operatorname{vol}(B^3) - \operatorname{vol}(P_N)$.
We present numerical evidence for the exact value of the variance constant:
\[
  \operatorname{Var}(\Delta_N) \;=\; \frac{576\pi}{N^3} + o(1/N^3).
\]
Normalizing by $\operatorname{vol}(B^3) = 4\pi/3$, the variance constant
becomes $324/\pi$.  This appears not to have been computed in the literature,
despite extensive work on the asymptotic theory of random polytopes; it
answers a question posed by Akopyan, Edelsbrunner, and
Nikitenko~\cite{AEN2021}.  We also verify the central limit theorem for the
standardized deficit, with skewness and excess kurtosis both converging to
zero.
\end{abstract}

% ============================================================================
\section{Introduction}
% ============================================================================

The study of random polytopes---convex hulls of random point sets---is a
classical subject in stochastic geometry, with origins in the work of
R\'enyi and Sulanke~\cite{RS1963,RS1964}.  A basic question is: how
well does the convex hull of $N$ random points approximate the ambient
convex body as $N \to \infty$?

For the model considered here, $N$ independent points are sampled
uniformly on the unit sphere $S^{d-1} \subset \mathbb{R}^d$, and
$P_N$ denotes their convex hull.  Since every point lies on $S^{d-1}$,
every point is a vertex of $P_N$, and $P_N$ is an inscribed polytope
approximating the unit ball $B^d$ from within.  The volume deficit
$\Delta_N = \operatorname{vol}(B^d) - \operatorname{vol}(P_N)$
measures the quality of this approximation.

The expected volume of $P_N$ for $d = 3$ is known exactly.
Buchta, M\"uller, and Tichy~\cite{BMT1985} computed the expected
intrinsic volumes of random polytopes inscribed in spherical
bodies; Kabluchko, Temesvari, and Thal\"e~\cite{KTT2019} obtained
equivalent results for beta polytopes.  Akopyan, Edelsbrunner,
and Nikitenko~\cite{AEN2021} gave an elegant direct proof and
highlighted the particularly clean form for~$S^2$:
\begin{equation}\label{eq:exact-mean}
  \mathbb{E}[\operatorname{vol}(P_N)]
  = \frac{4\pi}{3}\,
    \frac{(N-1)(N-2)(N-3)}{(N+1)(N+2)(N+3)},
\end{equation}
from which the volume deficit satisfies
\begin{equation}\label{eq:deficit-exact}
  \mathbb{E}[\Delta_N]
  = \frac{4\pi}{3}\cdot
    \frac{12N^2 + 12}{(N+1)(N+2)(N+3)}
  = \frac{16\pi}{N} - \frac{96\pi}{N^2} + O(1/N^3).
\end{equation}

The asymptotic behavior of $\Delta_N$ has been studied extensively
beyond the mean.  For smooth convex bodies, B\'ar\'any~\cite{Barany2008}
surveys the general theory; Reitzner~\cite{Reitzner2005} establishes
central limit theorems for the volume of random polytopes;
Richardson, Vu, and Wu~\cite{RVW2007,RVW2008} treat the inscribed
case (vertices on the boundary), proving variance bounds
$c_1/N^3 \le \operatorname{Var}(\Delta_N) \le c_2/N^3$ for $d = 3$
and establishing a CLT; Calka, Schreiber, and Yukich~\cite{CSY2013}
establish Brownian limits for related functionals.

Despite this work, the \emph{explicit value} of the variance constant
\[
  \operatorname{Var}(\Delta_N) \sim \frac{C_2}{N^3}
\]
does not appear to have been computed for $d = 3$.  Akopyan,
Edelsbrunner, and Nikitenko~\cite{AEN2021} explicitly pose the
question of determining the \emph{distributions} of the normalized
intrinsic volume deficiencies (their Open Question~5).  In this
note, we answer part of that question by determining $C_2$
numerically to high precision and conjecturing its exact value.

% ============================================================================
\section{Experimental Setup}
% ============================================================================

We sample $N$ independent uniform points on $S^2$ by drawing
three-dimensional standard Gaussian vectors and normalizing to unit
length.  For each sample, we compute the convex hull volume using the
Quickhull algorithm (via \texttt{scipy.spatial.ConvexHull}).  The
deficit $\Delta_N = \operatorname{vol}(B^3) -
\operatorname{vol}(P_N)$ is recorded for each trial.

For each value of $N$, we compute $M$ independent trials and record
the sample statistics of the rescaled deficit $N \cdot \Delta_N$ and
the variance-rescaled quantity $N^{3/2} \cdot \Delta_N$.  Specifically,
we compute:
\begin{itemize}
  \item The sample mean of $N \cdot \Delta_N$ (estimating $C_1$),
  \item $N^{3/2}$ times the sample standard deviation of $\Delta_N$
    (estimating $\sqrt{C_2}$),
  \item The sample skewness and excess kurtosis of $N \cdot \Delta_N$
    (testing the CLT).
\end{itemize}

We use $M = 500{,}000$ trials for $N \le 10{,}000$ and $M =
50{,}000$--$200{,}000$ for larger $N$ (up to $N = 50{,}000$).
Computation is parallelized over 64 CPU cores.

% ============================================================================
\section{Results}
% ============================================================================

\subsection{The mean: verification of the exact formula}

Table~\ref{tab:data} reports the rescaled mean deficit $\hat{C}_1(N)
= N \cdot \overline{\Delta_N}$ for each $N$.  The sequence increases
monotonically toward $16\pi = 50.26548\ldots$, matching the
leading term of the exact formula~\eqref{eq:deficit-exact}.  In
fact, the agreement extends beyond the leading term: for each~$N$,
the sample mean of $\Delta_N$ agrees with the exact rational
expression~\eqref{eq:deficit-exact} to within statistical error,
providing a strong consistency check on our computations.

The exact formula also yields the full asymptotic expansion.
In terms of the fractional deficit,
\[
  \mathbb{E}\!\left[\frac{\Delta_N}{\operatorname{vol}(B^3)}\right]
  = \frac{12}{N} - \frac{72}{N^2} + \frac{312}{N^3} - \cdots,
\]
where the leading normalized constant $12$ is particularly striking
in its simplicity.

\subsection{The variance constant}

We estimate $\sqrt{C_2}$ via $N^{3/2} \cdot
\operatorname{sd}(\Delta_N)$.  Table~\ref{tab:data} shows this
quantity converging to approximately $42.54$.  We observe that
$24\sqrt{\pi} = 42.5389\ldots$, so that $C_2 = 576\pi$.

\begin{conjecture}\label{conj:var}
Under the same hypotheses,
\[
  \operatorname{Var}\!\left[\operatorname{vol}(B^3) -
  \operatorname{vol}(P_N)\right] = \frac{576\pi}{N^3} + o(1/N^3).
\]
Equivalently, the variance of the fractional deficit is $324/(\pi
N^3) + o(1/N^3)$.
\end{conjecture}

\subsection{Central limit theorem}

The skewness and excess kurtosis of $N \cdot \Delta_N$ both
decrease toward zero as $N$ increases (Table~\ref{tab:data}),
consistent with the CLT established by Richardson, Vu, and
Wu~\cite{RVW2008}.  Combining the exact mean with
Conjecture~\ref{conj:var}, the standardized deficit
\[
  Z_N = \frac{\Delta_N - \mathbb{E}[\Delta_N]}{24\sqrt{\pi}\,/\,N^{3/2}}
\]
converges in distribution to $\mathcal{N}(0,1)$.

\begin{table}[t]
\centering
\caption{Numerical estimates of the deficit constants.  $\hat{C}_1(N) = N \cdot \overline{\Delta_N}$ estimates the mean constant;
$\hat{\sigma}(N) = N^{3/2} \cdot \operatorname{sd}(\Delta_N)$ estimates $\sqrt{C_2}$.  Reference values: $16\pi \approx 50.2655$, $24\sqrt{\pi} \approx 42.5389$.}
\label{tab:data}
\smallskip
\begin{tabular}{@{}rcccc@{}}
\toprule
$N$ & $\hat{C}_1(N)$ & $\hat{\sigma}(N)$ & Skewness & Ex.\ Kurtosis \\
\midrule
10      & 29.589 &  8.646 & 0.363 & $-0.109$ \\
20      & 37.937 & 18.682 & 0.725 & 0.877 \\
50      & 44.723 & 30.454 & 0.795 & 1.334 \\
100     & 47.374 & 36.007 & 0.683 & 1.051 \\
200     & 48.790 & 39.220 & 0.531 & 0.643 \\
500     & 49.667 & 41.259 & 0.357 & 0.304 \\
1\,000  & 49.965 & 41.970 & 0.252 & 0.159 \\
2\,000  & 50.116 & 42.342 & 0.187 & 0.088 \\
5\,000  & 50.205 & 42.671 & 0.115 & 0.034 \\
10\,000 & 50.235 & 42.787 & 0.082 & 0.020 \\
20\,000 & 50.251 & 42.677 & 0.058 & 0.005 \\
30\,000 & 50.255 & 42.535 & 0.047 & 0.013 \\
50\,000 & 50.259 & 42.790 & 0.027 & 0.017 \\
\midrule
$\infty$ & $16\pi \approx 50.265$ & $24\sqrt{\pi} \approx 42.539$ & 0 & 0 \\
\bottomrule
\end{tabular}
\end{table}

% ============================================================================
\section{Discussion}
% ============================================================================

\subsection{Cap-volume heuristic for the variance}

The deficit $\Delta_N$ is the total volume of the $2N-4$ spherical
caps cut off by the faces of $P_N$.  A typical Delaunay triangle
has area $\sim 4\pi/N$, so the cap height is $h \sim 1/N$ and
the cap volume is $V_{\mathrm{cap}} \sim \pi/N^2$.  The variance
of a single cap volume is therefore $\operatorname{Var}(V_{\mathrm{cap}})
\sim c/N^4$.  Since cap volumes are only weakly dependent (caps
whose triangles are well-separated on the sphere are nearly
independent), summing $\sim 2N$ weakly dependent terms gives
\[
  \operatorname{Var}(\Delta_N) \sim 2N \cdot c/N^4 = O(1/N^3),
\]
consistent with the rigorous bounds of Richardson, Vu, and
Wu~\cite{RVW2008}.  However, extracting the exact constant
$576\pi$ from this heuristic would require the precise
distribution of cap volumes and their covariance structure.

\subsection{Relation to the literature}

The exact formula~\eqref{eq:exact-mean} for the expected volume
was first obtained (in more general form) by Buchta, M\"uller,
and Tichy~\cite{BMT1985} and later rederived by Kabluchko,
Temesvari, and Thal\"e~\cite{KTT2019} in the framework of beta
polytopes.  Akopyan, Edelsbrunner, and Nikitenko~\cite{AEN2021}
highlighted the clean form for~$S^2$ and gave an elegant
combinatorial proof using Crofton's formula.  They also posed
the question of determining the \emph{distributions} of the
normalized deficiencies (their Open Question~5); our variance
conjecture provides a partial answer.

Richardson, Vu, and Wu~\cite{RVW2008} prove that for
$d$-dimensional convex bodies with $C^2$ boundary, the variance
of the volume of the inscribed random polytope satisfies
$\operatorname{Var} \asymp N^{-(d+3)/(d-1)}$, which gives
$N^{-3}$ for $d = 3$.  Their methods involve Poissonization and
yield unspecified constants; the explicit value $576\pi$ does not
appear to follow from their framework.

\subsection{Higher dimensions}\label{sec:higher-d}

The exact mean formula generalizes to all~$d$: for $N$ uniform
points on $S^{d-1}$, the expected volume of the convex hull is
\[
  \mathbb{E}[\operatorname{vol}(P_N)]
  = \operatorname{vol}(B^d) \prod_{k=1}^{d} \frac{N-k}{N+k}.
\]
The normalized leading term of the fractional deficit is
$d(d+1)/N$, giving $12$ for $d = 3$, $20$ for $d = 4$, $30$
for $d = 5$, and $d(d+1)$ in general.  (Note: the pattern is
$d(d+1)$, not $2\cdot d!$ as might be guessed from the case
$d = 3$ alone.)

However, the deficit scaling changes qualitatively for $d \ge 4$.
For points on $S^{d-1}$ with $d \ge 4$, the expected deficit
scales as $N^{-2/(d-1)}$, not $N^{-1}$.  (The exponent
$2/(d-1)$ arises from the curvature of the
sphere~\cite{Muller1990}; for $d = 3$ one has $2/(d-1) = 1$,
recovering the $N^{-1}$ scaling.)  The product formula above
gives exact rational expressions for all~$N$, but the leading
asymptotics for $d = 4$ become
\[
  \mathbb{E}[\Delta_N] \sim C_d \cdot N^{-2/3},
\]
with convergence to the asymptotic regime that is much slower
than in $d = 3$: the correction terms are $O(N^{-1/3})$
relative to the leading term, compared to $O(N^{-1})$ for
$d = 3$.

We performed experiments for $d = 4$ with $N$ up to
$100{,}000$.  The rescaled deficit $N^{2/3}\Delta_N$ converges
very slowly, reaching $47.53$ at $N = 100{,}000$ (between the
candidates $15\pi \approx 47.12$ and $48$).  The variance
exponent, which Richardson, Vu, and Wu~\cite{RVW2008} predict
to be $(d+3)/(d-1) = 7/3$ for $d = 4$, is confirmed by our
data: the local log-log slopes of $\operatorname{Var}(\Delta_N)$
converge to $2.33$ by $N = 20{,}000$.  However, the variance
\emph{constant} has not yet converged at $N = 100{,}000$,
reflecting the same slow-convergence phenomenon.
Table~\ref{tab:data4d} summarizes the $d = 4$ results.

\begin{table}[t]
\centering
\caption{Volume deficit statistics for $d = 4$.
$N^{2/3}\overline{\Delta_N}$ estimates the mean constant;
$N^{7/3}\widehat{\operatorname{Var}}$ tests the RVW variance
scaling.  The log-log slope of the variance converges to the
predicted $7/3 \approx 2.333$.}
\label{tab:data4d}
\smallskip
\begin{tabular}{@{}rccc@{}}
\toprule
$N$ & $N^{2/3}\overline{\Delta_N}$ & $N^{7/3}\widehat{\operatorname{Var}}$ & Var.\ slope \\
\midrule
100     & 39.96 &  167.5 & --- \\
200     & 42.67 &  236.9 & 1.83 \\
500     & 44.89 &  305.5 & 2.06 \\
1\,000  & 45.89 &  340.2 & 2.18 \\
2\,000  & 46.52 &  363.2 & 2.24 \\
5\,000  & 47.02 &  380.2 & 2.28 \\
10\,000 & 47.24 &  390.3 & 2.30 \\
20\,000 & 47.38 &  391.7 & 2.33 \\
50\,000 & 47.48 &  399.4 & 2.31 \\
100\,000& 47.53 &  408.3 & 2.30 \\
\bottomrule
\end{tabular}
\end{table}

\subsection{Open questions}

\begin{enumerate}
  \item \emph{Prove} Conjecture~\ref{conj:var}.  The variance
    computation requires understanding the covariance structure
    of cap volumes across faces of the Delaunay triangulation,
    which goes beyond the techniques used for the mean
    in~\cite{AEN2021}.
  \item Compute the analogous variance constant for $d = 4$.
    Our experiments confirm the RVW exponent $7/3$ but do not
    yet pin down the constant.  Is there a clean closed form
    analogous to $576\pi$?
  \item Characterize the rate of convergence in the CLT.  Our
    data for $d = 3$ show persistent positive skewness
    $\sim N^{-1/2}$, suggesting that a Berry--Esseen bound
    of order $N^{-1/2}$ may be sharp.
  \item Determine the full distribution of the normalized
    deficiency, answering the remainder of Open Question~5
    of~\cite{AEN2021}.
\end{enumerate}

% ============================================================================
\section{Ideal polytope volumes in hyperbolic 4-space}\label{sec:hyp4d}
% ============================================================================

We now turn from Euclidean inscribed polytopes to ideal polytopes
in hyperbolic space.  An \emph{ideal polytope} in $\mathbb{H}^4$
is the convex hull of $N$ points on $S^3 = \partial\mathbb{H}^4$.
Its hyperbolic volume can be computed by decomposing into ideal
4-simplices via Delaunay triangulation and summing their volumes
using the Hopf--Poincar\'e--Schl\"afli angle formula.

\subsection{Volume of an ideal 4-simplex}

For a single ideal 4-simplex $[\infty, v_1, v_2, v_3, v_4]$ in
the upper half-space model, the volume is an \emph{exact affine
function} of the ten codimension-2 dihedral angles:
\begin{equation}\label{eq:vol4simplex}
  \operatorname{Vol}_4(\Delta)
  = \frac{\pi}{3}\Bigl(4\pi - \sum_{k=1}^{10}\theta_k\Bigr).
\end{equation}
The Euclidean tetrahedron $T = \operatorname{conv}(v_1,\ldots,v_4)$
determines all ten angles: six are the Euclidean dihedral angles
$\phi_{ij}$ of~$T$, and four are $\arccos(d_\ell/R)$, where $R$
is the circumradius and $d_\ell$ is the signed distance from the
circumcenter to the face opposite~$v_\ell$.

Because formula~\eqref{eq:vol4simplex} involves only $\arccos$
of algebraic numbers, the volume can be evaluated to \emph{arbitrary
precision} using extended arithmetic---no numerical integration
is needed.

\begin{remark}
The regular ideal 4-simplex has all ten dihedral angles equal to
$\arccos(1/3)$, giving
$\operatorname{Vol}_4(\Delta_{\mathrm{reg}})
= \frac{\pi}{3}(4\pi - 10\arccos(1/3)) \approx 0.26890$.
\end{remark}

\subsection{Volume distribution for random ideal 4-polytopes}

For $N$ uniformly random points on $S^3$, we computed the total
hyperbolic 4-volume of the corresponding ideal polytope
($100{,}000$ samples per vertex count, $N = 6$ to $100$).
Key findings:

\begin{enumerate}
\item The expected volume grows as $\mathbb{E}[V] \approx cN -
  b\sqrt{N}$, with $c \approx 1.10$ and a substantial $\sqrt{N}$
  correction.

\item A beta distribution fits the marginal volume distribution
  better than a normal (KS statistics $0.002$--$0.005$ for
  $N \ge 15$), mirroring the 3D behavior.  However, the beta
  parameters $\alpha$, $\beta$ scale as $N^{1.40}$ and $N^{1.20}$
  respectively, without the clean interpretation found in 3D.

\item The explanation lies in the \emph{fluctuating Delaunay
  combinatorics}.  In 3D, the number of Delaunay triangles for
  $N$ points on~$S^2$ is essentially deterministic ($2N - 4$
  minus boundary).  In 4D, the number of Delaunay tetrahedra for
  $N$ points on~$S^3$ has substantial variance: $\mathrm{CV}
  \approx 13\%$ at $N = 10$, still $1.7\%$ at $N = 100$.

\item The volume decomposes as $V = n \cdot \bar{v}$, where $n$
  is the number of Delaunay tetrahedra and $\bar{v}$ is the mean
  per-tetrahedron volume.  Both $\bar{v}$ and $n$ are
  approximately normally distributed, but they are \emph{negatively
  correlated} ($\rho \approx -0.18$ at $N = 100$): denser
  triangulations produce more but thinner simplices.

\item Conditioned on $n$, the volume is well-described by a normal
  distribution (KS~$\approx 0.02$), while the beta fit fails.
  The per-tetrahedron volume $\bar{v} = V/n$ is also normally
  distributed.  The marginal beta-like shape arises from the
  product of two negatively correlated normals, not from any
  intrinsic beta structure.
\end{enumerate}

\subsection{Optimal ideal 4-polytopes}

For each $N$, we maximized the total hyperbolic volume over all
configurations of $N$ points on~$S^3$, using L-BFGS-B with
200 random restarts and two-phase polishing.
Table~\ref{tab:opt4d} summarizes the results.

\begin{table}[t]
\centering
\caption{Optimal ideal 4-polytopes.  The $f$-vector
$(f_0, f_1, f_2, f_3)$ counts vertices, edges, triangles, and
tetrahedra of the boundary.  $|\mathrm{Aut}|$ is the automorphism
group order of the 1-skeleton (meaningful only when the skeleton
is not the complete graph~$K_N$, i.e., $N \ge 7$).}
\label{tab:opt4d}
\smallskip
\begin{tabular}{@{}rrcrl@{}}
\toprule
$N$ & Volume & $f$-vector & $|\mathrm{Aut}|$ & Identification \\
\midrule
 5  & 0.26890 & $(5,10,10,5)$      & 120 & Regular 4-simplex \\
 6  & 0.58783 & $(6,15,18,9)$      & 720$^*$ & Neighborly ($K_6$) \\
 7  & 1.02360 & $(7,19,24,12)$     &  48 & \\
 8  & 1.77143 & $(8,24,32,16)$     & 384 & 16-cell (hyperoctahedron) \\
 9  & 2.31385 & $(9,29,40,20)$     &  80 & \\
10  & 3.00910 & $(10,35,50,25)$    & 200 & \\
12  & 4.37720 & $(12,46,68,34)$    &   8 & \\
16  & 7.79862 & $(16,69,106,53)$   &   2 & \\
20  & 11.9358 & $(20,92,144,72)$   &  16 & \\
30  & 22.0037 & $(30,154,248,124)$ &   1 & \\
\bottomrule
\end{tabular}
\smallskip

{\footnotesize $^*$For $N \le 6$, points on $S^3$ in general
position form a neighborly polytope with 1-skeleton~$K_N$, so
$|\mathrm{Aut}(K_N)| = N!$ is not informative about geometric
symmetry.}
\end{table}

The optimizer for $N = 8$ is the \emph{16-cell}
(4-cross-polytope), with $f$-vector $(8,24,32,16)$ and
$|\mathrm{Aut}| = 384 = 2^4 \cdot 4!$ (the hyperoctahedral
group~$B_4$).  Using the exact coordinates $\pm e_i \in S^3$
and stereographic projection, each of the 8~congruent ideal
4-simplices has volume computable to arbitrary precision:
\[
  V_{\mathrm{tet}} = \frac{\pi}{3}\Bigl(\frac{3\pi}{2}
  - 6\arccos\frac{1}{\sqrt{3}} + \arccos\frac{1}{3}\Bigr),
  \qquad
  V_{\mathrm{16\text{-}cell}} = 8\,V_{\mathrm{tet}}
  \approx 1.77143.
\]
The extreme symmetry of the 16-cell also admits a direct
computation in the Klein model: the hyperbolic volume
element is $dV = d^4x\,(1 - |x|^2)^{-5/2}$, and the
16-cell is the $L^1$~unit ball $\{x : \sum|x_i| \le 1\}$,
so the volume is $16$ times a simplex integral.  This yields
the equivalent closed form
\begin{equation}\label{eq:16cell}
  V_{\mathrm{16\text{-}cell}}
  = \frac{4\pi^2}{3} - \frac{32\pi}{3}\,\arcsin\frac{1}{3}
  = \frac{4\pi^2}{15} - \frac{16}{5}\,v_4,
\end{equation}
where $v_4 = \frac{10\pi}{3}\arcsin(1/3) - \frac{\pi^2}{3}$
is the volume of the regular ideal 4-simplex.

In contrast, the optimizer for $N = 9$ has $|\mathrm{Aut}| = 80$
with two vertex orbits (of sizes 4 and~5), degree sequence
$[7^4, 6^5]$, and 10 congruent ideal 4-simplices involving
exactly 5~distinct angle values (including $2\pi/5$ and $\pi/2$).
However, all $\binom{9}{2} = 36$ pairwise inner products on~$S^3$
are distinct: the graph automorphism group is purely combinatorial,
not a metric symmetry.  The two non-rational angles
$\theta_1 \approx 51.03°$ and $\theta_2 \approx 62.81°$ satisfy
$\cos(2\theta_1) = 2\cos^2\theta_1 - 1$ (accounting for the
fifth angle type at $102.05°$), but do not appear to be
expressible in elementary closed form.  No Klein-model reduction
is available, and the volume $V(9) \approx 2.31385$ is accessible
only through the angle formula or numerical optimization.

\paragraph{Comparison with symmetry-based constructions.}
An appealing strategy for constructing large-volume ideal
4-polytopes is to exploit the supplementarity
$\arccos(1/3) + \arccos(-1/3) = \pi$ of the dihedral angles of
the regular ideal 4-simplex and the ideal 16-cell: gluing a
regular simplex onto one facet of the 16-cell produces a convex
ideal polytope with 9 vertices and volume
$V_{\mathrm{16\text{-}cell}} + v_4 \approx 2.040$.  While this
exceeds the mean volume of a random 9-vertex configuration
($\mathbb{E}[V] \approx 1.08$), it falls $12\%$ short of the
actual optimum $V(9) \approx 2.314$.  The optimizer achieves its
advantage not through regularity of the individual simplices but
through packing efficiency: 10~simplices at
$\approx 0.231$ each, versus $9$~simplices (8~at
$\approx 0.221$ plus 1~at $0.269$) in the glued construction.
This illustrates a general principle: \emph{maximal symmetry does
not imply maximal volume}, and the gap between symmetry-based
lower bounds and actual optima grows with~$N$.

\paragraph{Surface area and isoperimetric ratio.}
Each boundary facet of an ideal 4-polytope is an ideal tetrahedron
in a totally geodesic $\mathbb{H}^3 \subset \mathbb{H}^4$; its
hyperbolic 3-volume can be computed by the classical Lobachevsky
formula.  The total 3-surface area
$A = \sum_{\text{facets}} \operatorname{Vol}_3(F)$ defines a
natural companion to the 4-volume.
Table~\ref{tab:surface} compares the maximizers of volume, surface
area, and the isoperimetric ratio $V/A$ for small~$N$.

\begin{table}[t]
\centering
\caption{Three optimization objectives for ideal 4-polytopes.
For $N \le 10$ (excluding $N = 6$), all three objectives are
maximized by the \emph{same} polytope.  At $N = 6$ and
$N \ge 12$ the optimizers diverge.}
\label{tab:surface}
\smallskip
\begin{tabular}{@{}r|ccc|ccc|cc@{}}
\toprule
& \multicolumn{3}{c|}{Max Volume} & \multicolumn{3}{c|}{Max Area} & \multicolumn{2}{c}{Max $V/A$} \\
$N$ & $V$ & $A$ & $V/A$ & $A$ & $V$ & $V/A$ & $V/A$ & same? \\
\midrule
5   & 0.269 & 5.074 & 0.053 & 5.075 & 0.269 & 0.053 & 0.053 & yes \\
6   & 0.588 & 7.818 & 0.075 & 8.119 & 0.538 & 0.066 & 0.075 & V=iso$\ne$A \\
7   & 1.024 & 11.83 & 0.087 & 11.83 & 1.024 & 0.087 & 0.087 & yes \\
8   & 1.771 & 16.24 & 0.109 & 16.24 & 1.771 & 0.109 & 0.109 & yes \\
9   & 2.314 & 19.96 & 0.116 & 19.96 & 2.314 & 0.116 & 0.116 & yes \\
10  & 3.009 & 23.91 & 0.126 & 23.91 & 3.009 & 0.126 & 0.126 & yes \\
12  & 4.377 & 32.33 & 0.135 & 32.43 & 4.015 & 0.124 & 0.136 & all differ \\
16  & 7.799 & 50.49 & 0.154 & 50.74 & 7.785 & 0.153 & 0.155 & all differ \\
20  & 11.86 & 70.35 & 0.169 & 70.58 & 11.81 & 0.167 & 0.168 & all differ \\
\bottomrule
\end{tabular}
\end{table}

A striking finding: for $N \in \{5, 7, 8, 9, 10\}$, the same
polytope simultaneously maximizes all three objectives (the
entries marked ``yes'' in Table~\ref{tab:surface}).  These are
precisely the highly symmetric optimizers---the regular simplex
($N = 5$), the 16-cell ($N = 8$), and the configurations with
$|\mathrm{Aut}| \ge 48$.  At $N = 6$, the surface-area maximizer
has 8~boundary facets all at maximum volume~$v_3$ (the regular
ideal tetrahedron in $\mathbb{H}^3$), while the volume maximizer
has 9~facets at sub-maximal individual volume.  For $N \ge 12$,
all three objectives diverge, with the surface-area maximizer
favoring fewer but more regular boundary tetrahedra.

\paragraph{Dihedral angles and rationality.}
For each optimizer, we examined whether the dihedral angles are
rational multiples of~$\pi$ (which would be needed to tile
$\mathbb{H}^4$ by copies of the polytope and obtain a hyperbolic
manifold).  The cosines of the dihedral angles are
\emph{algebraic numbers of low degree}---square roots of
rationals---but the angles themselves are generically
\emph{not} rational multiples of~$\pi$.  In detail:
\begin{itemize}
\item $N=5$: all angles $\arccos(1/3)$;
\item $N=6$: angles $\arccos(\sqrt{3/7})$, $\arccos(-1/7)$, $2\pi/3$;
\item $N=8$: angles $\arccos(1/\!\sqrt{3})$, $\pi/2$, $\arccos(-1/3)$;
\item $N=9$: includes $2\pi/5$ and $\pi/2$ alongside irrational-over-$\pi$ angles.
\end{itemize}
Only $\pi/2$, $2\pi/3$, and $2\pi/5$ appear as rational multiples
of~$\pi$; the dominant angles are irrational over~$\pi$ in every
case.  This means these optimal ideal 4-polytopes cannot directly
tile $\mathbb{H}^4$ to produce manifolds.

A sharp symmetry phase transition occurs around $N = 12$:
for $N \le 10$, the optimizers have large automorphism groups
($48$--$384$, excluding the trivial $K_N$ cases), few distinct
angle values ($1$--$8$), and all tetrahedra are nearly congruent.
For $N \ge 12$, symmetry collapses to $|\mathrm{Aut}| \le 16$,
with hundreds of distinct angles and irregular structure.

\subsection{Facet-parameterized optimization}

The vertex-based optimization of Section~\ref{tab:opt4d}
parameterizes ideal 4-polytopes by their vertices on~$S^3$
and fixes the vertex count~$N$.  A complementary approach
parameterizes by \emph{facet normals}: given $f$~unit
vectors $u_1, \ldots, u_f \in S^3$, define
\[
  P(h) = \bigl\{ x \in \mathbb{R}^4 : u_i \cdot x \le h_i
  \;\text{for all } i \bigr\}.
\]
This polytope has at most~$f$ facets, with the number of vertices
determined by the arrangement of the bounding hyperplanes.
To obtain an \emph{ideal} polytope, we require all vertices of~$P(h)$
to lie on~$S^3$, i.e., $|v|^2 = 1$ for each vertex~$v$.

For a fixed combinatorial type (determined by the normals~$u_i$
at a probe value of~$h$), each vertex is a linear function
$v_j = A_j h$ of the support numbers, where $A_j$ encodes the
four active hyperplane constraints at vertex~$j$.  The ideal
condition $|v_j|^2 = 1$ becomes a system of quadratic equations
in the~$f$ unknowns~$h_i$, which we solve by Levenberg--Marquardt
to machine precision.  This inner solver succeeds for $30$--$70\%$
of random normal configurations, with the success rate depending
on~$f$.

\paragraph{Two-level optimization.}
To maximize volume over the facet normals, we use a two-level
scheme: the outer loop optimizes the $3f$~hyperspherical angles
parameterizing the normals on~$S^3$ via Powell's method, while
the inner loop solves the ideal constraint for each candidate
set of normals.  Table~\ref{tab:facet_opt} summarizes the results.

\begin{table}[t]
\centering
\caption{Facet-parameterized optimal ideal 4-polytopes.
$f$~denotes the number of facet normals, $N$~the emergent vertex
count, and $V$~the hyperbolic 4-volume.  The scan column gives
the best volume from 100~random normal configurations (inner solve
only); the optimized column includes the outer Powell
optimization over normals.}
\label{tab:facet_opt}
\smallskip
\begin{tabular}{@{}rrrrrr@{}}
\toprule
$f$ & $N$ & $V_{\mathrm{scan}}$ & $V_{\mathrm{opt}}$
    & Improvement & $V/f^2$ \\
\midrule
 8  &  18 &  5.56 &   8.01 & 44\% & 0.125 \\
10  &  30 & 12.37 &  15.92 & 29\% & 0.159 \\
12  &  41 & 18.68 &  26.65 & 43\% & 0.185 \\
16  &  61 & 36.71 &  41.51 & 13\% & 0.162 \\
20  &  91 & 57.07 &  68.60 & 20\% & 0.172 \\
24  & 111 & 78.42 &  84.27 &  8\% & 0.146 \\
\bottomrule
\end{tabular}
\end{table}

Several features stand out.  First, the outer optimization
consistently improves volumes by $8$--$44\%$ over the best
random scan, demonstrating that the normal directions matter
significantly.  Second, the vertex count~$N$ is not fixed by~$f$
but emerges from the dual construction: the number of vertices
equals the number of facets of $Q = \operatorname{conv}(u_i/h_i)$,
which varies with the normal configuration.  For simplicial
4-polytopes, $N$~satisfies the Dehn--Sommerville relation
$f_3 = f_1 - f_0$, and the typical vertex count grows roughly
as $N \approx 4f$.

Third, the optimal volume scales approximately as
$V \approx 0.17\,f^2$.  The ratio $V/f^2$ in the last column is
nearly constant across the range $f = 8$ to~$24$, suggesting a
quadratic growth law.  This should be contrasted with the
vertex-parameterized scaling, which is closer to linear in~$N$
(cf.\ Table~\ref{tab:opt4d}).

The cross-validated volumes agree to machine precision with
those obtained from the independent general Schl\"afli formula
for hyperbolic 4-simplices~\cite{Murakami2011}, which handles
both compact and ideal vertices via Murakami's dilogarithm formula
for spherical link volumes.

% ============================================================================
\section*{Acknowledgments}
% ============================================================================

Computations were performed using \texttt{scipy.spatial.ConvexHull}
with parallelization over 64 CPU cores.  Code is available at
\url{https://github.com/igorrivin/ideal_poly_volume_toolkit}.

\begin{thebibliography}{10}

\bibitem{AEN2021}
A.~Akopyan, H.~Edelsbrunner, and A.~Nikitenko,
The beauty of random polytopes inscribed in the 2-sphere,
\emph{Experimental Mathematics} \textbf{34} (2025), 327--341.

\bibitem{Barany2008}
I.~B\'ar\'any,
Random polytopes, convex bodies, and approximation,
in \emph{Stochastic Geometry} (A.~Baddeley, I.~B\'ar\'any,
R.~Schneider, W.~Weil, eds.), Lecture Notes in Math.\ 1892,
Springer, 2007, pp.~77--118.

\bibitem{CSY2013}
P.~Calka, T.~Schreiber, and J.~E.~Yukich,
Brownian limits, local limits and variance asymptotics for convex
hulls in the ball,
\emph{Ann.\ Probab.}\ \textbf{41} (2013), 50--108.

\bibitem{BMT1985}
C.~Buchta, J.~M\"uller, and R.~F.~Tichy,
Stochastical approximation of convex bodies,
\emph{Math.\ Ann.}\ \textbf{271} (1985), 225--235.

\bibitem{KTT2019}
Z.~Kabluchko, D.~Temesvari, and C.~Thal\"e,
Expected intrinsic volumes and facet numbers of random
beta-polytopes,
\emph{Math.\ Nachr.}\ \textbf{292} (2019), 79--105.

\bibitem{Muller1990}
J.~S.~M\"uller,
Approximation of a ball by random polytopes,
\emph{J.\ Approx.\ Theory} \textbf{63} (1990), 198--209.

\bibitem{Reitzner2005}
M.~Reitzner,
Central limit theorems for random polytopes,
\emph{Probab.\ Theory Related Fields} \textbf{133} (2005),
483--507.

\bibitem{RS1963}
A.~R\'enyi and R.~Sulanke,
\"Uber die konvexe H\"ulle von $n$ zuf\"allig gew\"ahlten Punkten,
\emph{Z.\ Wahrsch.\ Verw.\ Gebiete} \textbf{2} (1963), 75--84.

\bibitem{RS1964}
A.~R\'enyi and R.~Sulanke,
\"Uber die konvexe H\"ulle von $n$ zuf\"allig gew\"ahlten Punkten, II,
\emph{Z.\ Wahrsch.\ Verw.\ Gebiete} \textbf{3} (1964), 138--147.

\bibitem{RVW2007}
C.~Richardson, V.~H.~Vu, and L.~Wu,
Random inscribing polytopes,
\emph{European J.\ Combin.}\ \textbf{28} (2007), 2057--2071.

\bibitem{RVW2008}
C.~Richardson, V.~H.~Vu, and L.~Wu,
An inscribing model for random polytopes,
\emph{Discrete Comput.\ Geom.}\ \textbf{39} (2008), 469--499.

\bibitem{Murakami2011}
J.~Murakami,
The volume formulas for a spherical tetrahedron,
\emph{Proc.\ Amer.\ Math.\ Soc.}\ \textbf{140} (2012), 3289--3295.

\bibitem{Kolpakov2014}
A.~Kolpakov,
On the optimality of the ideal right-angled 24-cell,
\emph{Algebr.\ Geom.\ Topol.}\ \textbf{12} (2012), 1941--1960.

\end{thebibliography}

\end{document}
